% Part: proof-theory
% Chapter: proof-search
% Section: search-algorithm

\documentclass[../../../include/open-logic-section]{subfiles}

\begin{document}

\olfileid{pt}{ps}{sal}

\olsection{\Log{G3c} 的搜索算法}

本节描述 \Log{G3c} 的一个搜索算法。它并非唯一可能的算法，甚至也不是最有效的算法。
但它是正确的：如果一个相继式 $\Gamma \Sequent \Delta$ 有推导，该算法最终将终止并找到一个推导。
它也是万无一失的：如果算法终止却没有找到推导，或者永不终止，那么 $\Gamma
\Sequent \Delta$ 从一开始就不是有效的。

该算法一个关键的必要特征是，它确保 $\Gamma \Sequent \Delta$ 中、或证明搜索过程中产生的任一相继式中的每个公式都被“归约”，
即，它被作为反向应用的推理规则的主公式来处理，并且这一过程会根据需要重复足够多次。
我们称这一特性为\emph{公平性}。

该算法分阶段运行。每一阶段的输出是一个相继式树，
下一阶段的输出则是通过归约每个尚未是公理的顶端相继式中的一个公式而产生的。

在第 $0$ 阶段，我们从 $\Gamma \Sequent \Delta$ 开始。
为确保公平性，我们给 $\Gamma$ 和 $\Delta$ 中的公式分配数字作为指标，使得不同公式获得不同的编号。
比如，可以设想从左向右为公式编号。
我们将编号写为上标。
例如，若要搜索 $!A, !B \Sequent !C$ 的一个推导，那么第$0$阶段的输出是 $!A^1, !B^2 \Sequent !C^3$。
在这样的指标相继式中，作为上标出现的最大~$n$ 称为它的\emph{最大指标}（示例中为~$3$）。

如果相继式包含量词，我们还需知道当归约右侧的 $\lexists[x][!A(x)]$ 或左侧的 $\lforall[x][!A(x)]$ 时该使用哪些项。
我们只需由常项符号 $\Obj c_0$, $\Obj c_1$, $\Obj c_2$, \dots 及出现在~$\Gamma$ 和 $\Delta$ 中的函数符号所构造出的项。
我们假设所有这些项的集合已按某种固定的枚举给出，这样我们就可以谈论诸如“第一个不在 $\Gamma \Sequent \Delta$ 中出现的常项符号”之类的概念。
算法生成的树中的每个指标相继式都有一个关联的常项符号集合。
在第$0$阶段，指派给相继式的常项符号集合是其中出现的所有常项符号的集合（如果有的话），如果它不含任何常项符号，则是 $\{\Obj c_0\}$。

假设算法在第~$n$ 阶段已经生成了一棵带有关联常项符号集合的指标相继式树。
在第~$n+1$ 阶段，我们对每个顶端相继式执行以下操作。

如果有一个公式~$!A$ 同时出现在相继式的左边和右边，或者左边包含~$\lfalse$，我们不做任何操作：
这样的相继式在~\Log{G3c} 中已有推导，对于该顶端相继式无需再作处理。

如果相继式中没有非原子公式，则中止算法。$\Gamma \Sequent \Delta$ 没有证明。

否则，选取指标~$i$ 最小的非原子公式~$!A$。那么该顶端相继式要么是 $!A^i, \Pi \Sequent \Lambda$，要么是 $\Pi \Sequent \Lambda, !A^i$。
令 $k$ 为该相继式的最大指标，$C$ 为指派给它的常项符号集合。

我们“归约” $!A^i$，即根据~$!A$ 的主联结词以及 $!A^i$ 出现在左边还是右边，按照相应的推理规则生成新的顶端相继式。

\begin{enumerate}
  \item 若 $!A \ident !B \land !C$ 且出现在左边：在 $(!B \land !C)^i, \Pi \Sequent \Lambda$ 上方添加一个新的顶端相继式：
  \[
  \Axiom$!B^k, !C^{k+1}, \Pi \fCenter \Lambda$
  \RightLabel{\LeftR{\land}}
  \UnaryInf$(!B \land !C)^i, \Pi \fCenter \Lambda$
  \DisplayProof
  \]
  指派给新相继式的常项符号集合是~$C$。
  \item 若 $!A \ident !B \land !C$ 且出现在右边：在 $\Pi \Sequent \Lambda, (!B \land !C)^i$ 上方添加两个新的顶端相继式：
  \[
  \Axiom$\Pi \fCenter \Lambda, !B^k$
  \Axiom$\Pi \fCenter \Lambda, !C^k$
  \RightLabel{\RightR{\land}}
  \BinaryInf$\Pi \fCenter \Lambda, (!B \land !C)^i$
  \DisplayProof
  \]
  指派给新相继式的常项符号集合是~$C$。

  \item $!A \ident \lnot !B$、$!A \ident !B \lor !C$ 以及 $!A \ident !B \lif !C$ 的情况类似，留作练习。

  \item 若 $!A \ident \lexists[x][!B(x)]$ 且出现在左边：在 $\lexists[x][!B(x)]^i, \Pi \Sequent \Lambda$ 上方添加一个新的顶端相继式：
  \[
  \Axiom$!B(c)^k, \Pi \fCenter \Lambda$
  \RightLabel{\LeftR{\lexists}}
  \UnaryInf$\lexists[x][!B(x)]^i, \Pi \fCenter \Lambda$
  \DisplayProof
  \]
  其中 $c$ 是不在~$C$ 中的第一个常项符号。指派给新相继式的常项符号集合是 $C \cup \{c\}$。
  \item 若 $!A \ident \lexists[x][!B(x)]$ 且出现在右边：在 $\Pi \Sequent \Lambda, \lexists[x][!B(x)]$ 上方添加一个新的顶端相继式：
  \[
  \Axiom$\Pi \fCenter \Lambda, \lexists[x][!B(x)]^k, !B(t)^{k+1}$
  \RightLabel{\LeftR{\lexists}}
  \UnaryInf$\Pi \fCenter \Lambda, \lexists[x][!B(x)]^i$
  \DisplayProof
  \]
  其中 $t$ 是仅由~$C$ 中常项符号构成、且在此相继式下方对 $\lexists[x][!B(x)]$ 进行右侧归约时尚未使用过的第一个项（若所有这样的项均已使用，则取 $\Obj c_0$）。指派给新相继式的常项符号集合是 $C$。
  \item $!A \ident \lforall[x][!B(x)]$ 的情况类似，留作练习。
\end{enumerate}

如果在阶段 $n+1$ 我们没有向任何顶端相继式添加新的相继式，那么算法成功终止。此时，我们通过擦除阶段 $n+1$ 的树中所有指标，便得到了~$\Gamma \Sequent \Delta$ 的一个推导。顶端相继式都具有 $!A, \Pi \Sequent \Lambda, !A$ 的形式，或具有 $\lfalse, \Pi \Sequent \Lambda$ 的形式。所有的推理都是~\Log{G3c} 的正确推理。这对命题推理而言是显然的。注意，在每一步，指派给一个相继式的常项符号集合~$C$ 包含了该相继式中出现的所有常项符号。因此，在使用 \LeftR{\lexists} 和 \RightR{\lforall} 进行归约时，本征变元条件是满足的，因为本征变元~$c$ 是一个不在结论的常项符号集合~$C$ 中的常项符号，所以 $c$ 也不出现在结论中。我们得到：

\begin{prop}
\Log{G3c} 的证明搜索算法是正确的：如果它成功终止，则起始相继式~$\Gamma \Sequent \Delta$ 有推导。
\end{prop}





\end{document}